% Original source preamble and source-level macro declarations.
% This archival file is not the blueprint driver preamble.

\documentclass[11pt, reqno]{amsart}

%----------------------------------------------Showkeys
%\usepackage[notcite]{showkeys}   % for drafts to show all the labels in pdf

%----------------------------------------------Colors and color comments!
\usepackage[usenames,dvipsnames]{color}
\newcommand{\kestutis}[1]{{\color{Plum} \sf \it{ {\scriptsize ``#1''} }}}
\newcommand{\ready}[1]{{\color{Brown}   
#1}}						% parts carefully proofread
\newcommand{\revise}[1]{{\color{Bittersweet}  
 #1}}					% written parts, not in final form
\newcommand{\temp}[1]{{\color{Gray} {\footnotesize  #1 }}}			% obsolete parts, leftovers
\newcommand{\me}[1]{{\color{Plum} \sf \it{ {\scriptsize ``#1''} }}}

\renewcommand{\ready}[1]{#1}	    	
\renewcommand{\revise}[1]{#1}   		
%\renewcommand{\temp}[1]{#1}     		

\newcommand{\coR}{\textcolor{red}}
\newcommand{\coG}{\textcolor{green}}
\newcommand{\coB}{\textcolor{blue}}

%----------------------------------------------Packages
\usepackage[OT2, T1]{fontenc}
\usepackage{url}
\usepackage{amsmath}
\usepackage{array}
\usepackage{graphicx}
%\usepackage{amsfonts}
\usepackage{amssymb}
\usepackage{amsthm}
\definecolor{link}{RGB}{11,0,128}
\usepackage[colorlinks=true,citecolor=NavyBlue,linkcolor=OliveGreen,urlcolor=link]{hyperref}
\usepackage{hyperref}
\usepackage{paralist}
\usepackage{colonequals}			% for nice :=
\usepackage{color}
\usepackage{mathabx}			% for \widec
\usepackage{amscd}
\usepackage[all,cmtip]{xy}			% for commutative diagrams
\usepackage{sseq}				% for spectral sequences
\usepackage{verbatim}
\usepackage{parskip}
\usepackage{microtype}
\usepackage[in]{fullpage}
%\usepackage{graphicx}
\usepackage{mathrsfs} 			% for \mathscr (script letters)
%\usepackage{bm} 				% for bold Greek letters
\usepackage{cleveref}
\usepackage{enumitem}
\usepackage{moreenum}			% for enumeration via Greek letters
\usepackage[alphabetic,lite]{amsrefs} 	% for bibliography
\usepackage{stmaryrd}			% this is for \llbracket and \rrbracket
%\usepackage[foot]{amsaddr}
\usepackage{subfiles}
\usepackage{ragged2e}
\usepackage[bbgreekl]{mathbbol}     % This package is in order to have the prism symbol
\DeclareSymbolFontAlphabet{\mathbb}{AMSb}   % This is in order to common blackboard letters in the usual style (else the prism package messes them up)
\usepackage{stackengine}


\setlength\marginparwidth{2cm}


%-----------------------Greek letters
\newcommand{\gA}{\alpha}
\newcommand{\gB}{\beta}
\newcommand{\gG}{\gamma}
\newcommand{\gD}{\delta}
\newcommand{\gE}{\epsilon}
\newcommand{\gZ}{\zeta}
\newcommand{\gI}{\iota}
\newcommand{\gK}{\kappa}
\newcommand{\gL}{\lambda}
\newcommand{\gM}{\mu}
\newcommand{\gN}{\nu}
\newcommand{\gX}{\xi}
\newcommand{\gC}{\chi}
\newcommand{\gO}{\omega}
\newcommand{\gP}{\pi}
\newcommand{\gF}{\phi}
\newcommand{\gT}{\tau}

%-----------------------Blackboard letters
\newcommand{\bA}{\mathbb{A}}
\newcommand{\bC}{\mathbb{C}}
\newcommand{\bD}{\mathbb{D}}
\newcommand{\bE}{\mathbb{E}}
\newcommand{\bF}{\mathbb{F}}
\newcommand{\bG}{\mathbb{G}}
\newcommand{\bH}{\mathbb{H}}
\newcommand{\bL}{\mathbb{L}}
\newcommand{\bM}{\mathbb{M}}
\newcommand{\bN}{\mathbb{N}}
\newcommand{\bP}{\mathbb{P}}
\newcommand{\bQ}{\mathbb{Q}}
\newcommand{\bR}{\mathbb{R}}
\newcommand{\bS}{\mathbb{S}}
\newcommand{\bT}{\mathbb{T}}
\newcommand{\bU}{\mathbb{U}}
\newcommand{\bV}{\mathbb{V}}
\newcommand{\bZ}{\mathbb{Z}}

%-----------------------Bold letters
\newcommand{\bbA}{\mathbf{A}}
\newcommand{\bbB}{\mathbf{B}}
\newcommand{\bbC}{\mathbf{C}}
\newcommand{\bbD}{\mathbf{D}}
\newcommand{\bbE}{\mathbf{E}}
\newcommand{\bbF}{\mathbf{F}}
\newcommand{\bbG}{\mathbf{G}}
\newcommand{\bbH}{\mathbf{H}}
\newcommand{\bbK}{\mathbf{K}}
\newcommand{\bbL}{\mathbf{L}}
\newcommand{\bbN}{\mathbf{N}}
\newcommand{\bbP}{\mathbf{P}}
\newcommand{\bbQ}{\mathbf{Q}}
\newcommand{\bbR}{\mathbf{R}}
\newcommand{\bbS}{\mathbf{S}}
\newcommand{\bbT}{\mathbf{T}}
\newcommand{\bbU}{\mathbf{U}}
\newcommand{\bbX}{\mathbf{X}}
\newcommand{\bbY}{\mathbf{Y}}
\newcommand{\bbZ}{\mathbf{Z}}

\newcommand{\bba}{\mathbf{a}}
\newcommand{\bbb}{\mathbf{b}}
\newcommand{\bbc}{\mathbf{c}}
\newcommand{\bbd}{\mathbf{d}}
\newcommand{\bbe}{\mathbf{e}}
\newcommand{\bbf}{\mathbf{f}}
\newcommand{\bbg}{\mathbf{g}}
\newcommand{\bbh}{\mathbf{h}}
\newcommand{\bbi}{\mathbf{i}}
\newcommand{\bbj}{\mathbf{j}}
\newcommand{\bbk}{\mathbf{k}}
\newcommand{\bbl}{\mathbf{l}}
\newcommand{\bbm}{\mathbf{m}}
\newcommand{\bbn}{\mathbf{n}}
\newcommand{\bbo}{\mathbf{o}}
\newcommand{\bbp}{\mathbf{p}}
\newcommand{\bbq}{\mathbf{q}}
\newcommand{\bbr}{\mathbf{r}}
\newcommand{\bbs}{\mathbf{s}}
\newcommand{\bbt}{\mathbf{t}}
\newcommand{\bbu}{\mathbf{u}}
\newcommand{\bbv}{\mathbf{v}}
\newcommand{\bbw}{\mathbf{w}}
\newcommand{\bbx}{\mathbf{x}}
\newcommand{\bby}{\mathbf{y}}
\newcommand{\bbz}{\mathbf{z}}

%-----------------------Calligraphic letters
\newcommand{\cA}{\mathcal{A}}
\newcommand{\cB}{\mathcal{B}}
\newcommand{\cC}{\mathcal{C}}
\newcommand{\cD}{\mathcal{D}}
\newcommand{\cE}{\mathcal{E}}
\newcommand{\cF}{\mathcal{F}}
\newcommand{\cG}{\mathcal{G}}
\newcommand{\cH}{\mathcal{H}}
\newcommand{\cI}{\mathcal{I}}
\newcommand{\cJ}{\mathcal{J}}
\newcommand{\cK}{\mathcal{K}}
\newcommand{\cL}{\mathcal{L}}
\newcommand{\cM}{\mathcal{M}}
\newcommand{\cN}{\mathcal{N}}
\newcommand{\cO}{\mathcal{O}}
\newcommand{\cP}{\mathcal{P}}
\newcommand{\cQ}{\mathcal{Q}}
\newcommand{\cR}{\mathcal{R}}
\newcommand{\cS}{\mathcal{S}}
\newcommand{\cT}{\mathcal{T}}
\newcommand{\cU}{\mathcal{U}}
\newcommand{\cV}{\mathcal{V}}
\newcommand{\cW}{\mathcal{W}}
\newcommand{\cX}{\mathcal{X}}
\newcommand{\cY}{\mathcal{Y}}
\newcommand{\cZ}{\mathcal{Z}}

%-----------------------Fraktur letters
\newcommand{\fa}{\mathfrak{a}}
\newcommand{\fb}{\mathfrak{b}}
\newcommand{\fc}{\mathfrak{c}}
\newcommand{\ff}{\mathfrak{f}}
\newcommand{\fk}{\mathfrak{k}}
\newcommand{\fl}{\mathfrak{l}}
\newcommand{\fm}{\mathfrak{m}}
\newcommand{\fn}{\mathfrak{n}}
\newcommand{\fo}{\mathfrak{o}}
\newcommand{\fp}{\mathfrak{p}}
\newcommand{\fq}{\mathfrak{q}}
\newcommand{\fr}{\mathfrak{r}}
\newcommand{\ft}{\mathfrak{t}}

\newcommand{\fA}{\mathfrak{A}}
\newcommand{\fB}{\mathfrak{B}}
\newcommand{\fC}{\mathfrak{C}}
\newcommand{\fD}{\mathfrak{D}}
\newcommand{\fE}{\mathfrak{E}}
\newcommand{\fF}{\mathfrak{F}}
\newcommand{\fG}{\mathfrak{G}}
\newcommand{\fH}{\mathfrak{H}}
\newcommand{\fI}{\mathfrak{I}}
\newcommand{\fK}{\mathfrak{K}}
\newcommand{\fL}{\mathfrak{L}}
\newcommand{\fM}{\mathfrak{M}}
\newcommand{\fN}{\mathfrak{N}}
\newcommand{\fO}{\mathfrak{O}}
\newcommand{\fP}{\mathfrak{P}}
\newcommand{\fQ}{\mathfrak{Q}}
\newcommand{\fR}{\mathfrak{R}}
\newcommand{\fS}{\mathfrak{S}}
\newcommand{\fT}{\mathfrak{T}}
\newcommand{\fU}{\mathfrak{U}}
\newcommand{\fV}{\mathfrak{V}}
\newcommand{\fW}{\mathfrak{W}}
\newcommand{\fX}{\mathfrak{X}}
\newcommand{\fY}{\mathfrak{Y}}
\newcommand{\fZ}{\mathfrak{Z}}

%-----------------------Script letters
\newcommand{\sA}{\mathscr{A}}
\newcommand{\sB}{\mathscr{B}}
\newcommand{\sC}{\mathscr{C}}
\newcommand{\sD}{\mathscr{D}}
\newcommand{\sE}{\mathscr{E}}
\newcommand{\sF}{\mathscr{F}}
\newcommand{\sG}{\mathscr{G}}
\newcommand{\sH}{\mathscr{H}}
\newcommand{\sI}{\mathscr{I}}
\newcommand{\sJ}{\mathscr{J}}
\newcommand{\sK}{\mathscr{K}}
\newcommand{\sL}{\mathscr{L}}
\newcommand{\sM}{\mathscr{M}}
\newcommand{\sN}{\mathscr{N}}
\newcommand{\sO}{\mathscr{O}}
\newcommand{\sP}{\mathscr{P}}
\newcommand{\sQ}{\mathscr{Q}}
\newcommand{\sR}{\mathscr{R}}
\newcommand{\sS}{\mathscr{S}}
\newcommand{\sT}{\mathscr{T}}
\newcommand{\sU}{\mathscr{U}}
\newcommand{\sV}{\mathscr{V}}
\newcommand{\sW}{\mathscr{W}}
\newcommand{\sX}{\mathscr{X}}
\newcommand{\sY}{\mathscr{Y}}
\newcommand{\sZ}{\mathscr{Z}}

%-----------------------New operators
\DeclareMathOperator{\Ani}{Ani}		% Animation
\DeclareMathOperator{\Ann}{Ann}		% Annihilator of a module
\DeclareMathOperator{\Ass}{Ass}		% The set of associated primes
\DeclareMathOperator{\Aut}{Aut}		% The group of automorphisms
\DeclareMathOperator{\Bl}{Bl}			% blowup
\DeclareMathOperator{\Br}{Br}			% The Brauer group
\DeclareMathOperator{\CH}{CH}		% Chow group
\DeclareMathOperator{\Cauchy}{Cauchy}	% The ring of Cauchy sequences
\DeclareMathOperator{\Char}{char}		% Characteristic of a field
\DeclareMathOperator{\cid}{cid}		% the complete intersection defect
\DeclareMathOperator{\Cl}{Cl}			% Class group
\DeclareMathOperator{\CM}{CM}		% Cohen--Macaulay locus
\DeclareMathOperator{\codim}{codim}	% codimension
\DeclareMathOperator{\Coeq}{Coeq}		% The coequalizer
\DeclareMathOperator{\Coker}{Coker}	% Cokernel
\DeclareMathOperator{\Cone}{Cone}		% The cone of a morphism
\DeclareMathOperator{\Cont}{Cont}		% The space of continuous valuations
\DeclareMathOperator{\cosk}{cosk}		% The coskeleton 
\DeclareMathOperator{\depth}{depth}	% Depth of a module
\DeclareMathOperator{\dist}{dist}		% Distance
\DeclareMathOperator{\edim}{edim}		% The embedding dimension
\DeclareMathOperator{\End}{End}		% The algebra of endomorphisms
\DeclareMathOperator{\Eq}{Eq}			% The equalizer
\DeclareMathOperator{\Ext}{Ext}		% Derived functors of Hom
\DeclareMathOperator{\Extrig}{Extrig}	% Rigidified extensions
\DeclareMathOperator{\Fib}{Fib}		% Mapping fiber
\DeclareMathOperator{\Fil}{Fil}			% filtration
\DeclareMathOperator{\Fitt}{Fitt}		% The Fitting ideal
\DeclareMathOperator{\Frac}{Frac}		% Field of fractions
\DeclareMathOperator{\Fun}{Fun}		% Homomorphisms between two modules, etc.
\DeclareMathOperator{\Gal}{Gal}		% Galois group
\DeclareMathOperator{\gdepth}{geom. depth}		% the geometric depth of a ring
\DeclareMathOperator{\GL}{GL}		% The general linear group
\DeclareMathOperator{\gr}{gr}			% The associated graded
\DeclareMathOperator{\HH}{HH}		% Hochschild homology
\DeclareMathOperator{\Hom}{Hom}		% Homomorphisms between two modules, etc.
\DeclareMathOperator{\HP}{HP}		% periodic cyclic homology
\DeclareMathOperator{\Ht}{ht}			% Height of an ideal
\DeclareMathOperator{\im}{Im}			% The image of a map
\DeclareMathOperator{\Ind}{Ind}		% The induced representation
\DeclareMathOperator{\inv}{inv}		% The invariant map from local class field theory
\DeclareMathOperator{\Irr}{Irr}			% The set of isomorphism classes of irreducible representations
\DeclareMathOperator{\Isom}{Isom}		% The functor of isomorphisms
\DeclareMathOperator{\Jac}{Jac}		% The Jacobian
\DeclareMathOperator{\Ker}{Ker}		% The kernel
\DeclareMathOperator{\LCI}{LCI}		% Local complete intersection locus
\DeclareMathOperator{\lcm}{lcm}		% The least common multiple
\DeclareMathOperator{\length}{length}	% The length of a module
\DeclareMathOperator{\Lie}{Lie}		% The Lie algebra
\DeclareMathOperator{\Mat}{Mat}		% Matrices
\DeclareMathOperator{\Mod}{Mod}		% Category of modules
\DeclareMathOperator{\Norm}{Norm}	% The norm
\DeclareMathOperator{\NS}{NS}		% The Neron--Severi group
\DeclareMathOperator{\Null}{Null}		% The nullspace or the ring of null sequences
\DeclareMathOperator{\Ob}{Ob}		% Objects of a category
\DeclareMathOperator{\ord}{ord}		% The order
\DeclareMathOperator{\pd}{pd}			% projective dimension
\DeclareMathOperator{\Perv}{Perv}		% The category of perverse sheaves
\DeclareMathOperator{\PGL}{PGL}		% Projective general linear group
\DeclareMathOperator{\Pic}{Pic}		% Picard group
\DeclareMathOperator{\Pro}{Pro}		% Pro category
\DeclareMathOperator{\Proj}{Proj}		% Proj of a graded ring
\DeclareMathOperator{\rad}{rad}		% Radical of an ideal
\DeclareMathOperator{\Reg}{Reg}		% regular locus
\DeclareMathOperator{\Res}{Res}		% The restriction of the representation
\DeclareMathOperator{\rk}{rk}			% rank
\DeclareMathOperator{\Sel}{Sel}		% Selmer group
\DeclareSymbolFont{cyrletters}{OT2}{wncyr}{m}{n}
\DeclareMathSymbol{\Sha}{\mathalpha}{cyrletters}{"58}	% Sha
\DeclareMathOperator{\sk}{sk}			% The skeleton
\DeclareMathOperator{\SL}{SL}			% Special linear group
\DeclareMathOperator{\Spa}{Spa}		% Adic spectrum of an f-adic ring
\DeclareMathOperator{\Spec}{Spec}		% Spectrum of a ring
\DeclareMathOperator{\Spf}{Spf}		% Formal spectrum of an adic ring
\DeclareMathOperator{\Sq}{Sq}			% Steenrod squares
\DeclareMathOperator{\Sw}{Sw}		% Swan conductor
\DeclareMathOperator{\Supp}{Supp}		% Support of a module
\DeclareMathOperator{\Sym}{Sym}		% symmetric power
\DeclareMathOperator{\TC}{TC}		% topological cyclic homology
\DeclareMathOperator{\THH}{THH}		% topological Hochschild homology
\DeclareMathOperator{\Tor}{Tor}		% Derived functors of tensor product
\DeclareMathOperator{\TP}{TP}		% periodic topological cyclic homology
\DeclareMathOperator{\Tr}{Tr}			% Trace
\DeclareMathOperator{\Trace}{Trace}	% Trace
\DeclareMathOperator{\trdeg}{tr.deg}		% The transcendence degree
\DeclareMathOperator{\vdim}{vdim}		% The virtual dimension
\DeclareMathOperator{\Vect}{Vect}		% Vector bundles
\DeclareMathOperator{\Ver}{Ver}		% Verschiebung


%-----------------------Mathematics shortcuts
\providecommand{\Ab}{\mathrm{Ab}}
\newcommand{\ab}{\mathrm{ab}}				% Abelianization
\newcommand{\ad}{\mathrm{ad}}			% associated adic space
\newcommand{\aff}{\mathrm{aff}}	% for affine
\newcommand{\alg}{\mathrm{alg}}		% algebraic closure (mainly used in superscripts)
\newcommand{\ani}{{\mathrm{ani}}}		% the anisotropic locus
\newcommand{\arc}{{\rm arc}}
\newcommand{\Art}{\mathrm{Art}}				% Artin homomorphism
\newcommand{\bal}{\mathrm{bal}}		% balanced level structures
\newcommand{\ce}{\colonequals}
\newcommand{\Cong}{\mathrm{cong}}	% Congruence number
\newcommand{\const}{\mathrm{const}}	% Constant
\newcommand{\cont}{{\mathrm{cont}}}		% Continuous 
\newcommand{\cris}{{\mathrm{cris}}}		% Crystalline
\newcommand{\cycl}{{\mathrm{cycl}}}	% cyclotomic
\newcommand{\del}{\partial}
\newcommand{\dirlim}{\varinjlim}
\newcommand{\Div}{{\mathrm{Div}}}		% Cartier divisor associated to a generically exact complex
\renewcommand{\div}{{\mathrm{div}}}	% Maximal divisible subgroup
\newcommand{\dR}{{\mathrm{dR}}}		% de Rham
\newcommand{\eps}{\epsilon}
\newcommand{\et}{\mathrm{\acute{e}t}}	% for etale cohomology (mainly used in subscripts)
\newcommand{\fet}{\mathrm{f\acute{e}t}}	% for etale cohomology (mainly used in subscripts)
\newcommand{\fppf}{\mathrm{fppf}}		% for fppf cohomology (mainly used in subscripts)
\newcommand{\fpqc}{\mathrm{fpqc}}		% for fppf cohomology (mainly used in subscripts)
\newcommand{\Frob}{\mathrm{Frob}}			% Frobenius
\providecommand{\Gp}{\mathrm{Gp}}
\newcommand{\gp}{{\mathrm{gp}}}		% Group completion of a monoid
\newcommand{\heart}{{\heartsuit}}		% Heart of a t-structure
\newcommand{\hra}{\hookrightarrow}
\DeclareMathOperator{\Idem}{\mathrm{Idem}}	% The set of idempotents
\renewcommand{\i}{^{-1}}
\newcommand{\id}{\mathrm{id}}			% Identity
\newcommand{\Inf}{\mathrm{inf}}		% for A_\Inf
\newcommand{\injects}{\hookrightarrow}
\newcommand{\invlim}{\varprojlim}
\newcommand{\isomto}{\overset{\sim}{\longrightarrow}}
\newcommand{\isomfrom}{\overset{\sim}{\longleftarrow}}
\newcommand{\La}{\Leftarrow}
\newcommand{\llb}{\llbracket}			% [[
\newcommand{\llp}{(\!(}			% ((
\newcommand{\loc}{\mathrm{loc}}			% Localization map
\newcommand{\MCM}{\underline{\mathrm{MCM}}}  % the stable cateogry of maximal CM modules
\newcommand{\naive}{\mathrm{naive}}		% was used for X_0(n)^\naive
\newcommand{\nd}{{\mathrm{nd}}}			% the maximal nondivisible quotient
\newcommand{\normal}{{\lhd}}			% normal subgroup
\newcommand{\nr}{{\mathrm{nr}}}			% unramified 
\newcommand{\nsp}{\mathrm{nsp}}		% nonsplit (for instance, or nonsplit Cartan curve)
\newcommand{\old}{\mathrm{old}}			% old quotient of a modular Jacobian
\newcommand{\op}{{\mathrm{op}}}			% Opposite category
\newcommand{\ov}{\overline}
\providecommand{\p}[1]{\left(#1\right)}
\newcommand{\PD}{{\mathrm{PD}}}		% Divided powers
\newcommand{\perf}{\mathrm{perf}}			% perfection
\newcommand{\pr}{^{\prime}}
\newcommand{\prpr}{^{\prime\prime}}
\DeclareSymbolFontAlphabet{\mathbbl}{bbold}
\newcommand{\Prism}{{ \mathbbl{\Delta}}}
\newcommand{\proet}{\mathrm{pro\acute{e}t}}			% proetale
\newcommand{\profppf}{\mathrm{pro\x{-}fppf}}		% pro-fppf site
\newcommand{\prosyn}{\mathrm{pro\x{-}syn}}		% pro-syn site
\newcommand{\psh}{\mathrm{psh}}		% presheaf
\newcommand{\pt}{\mathrm{pt}}			% A point
\newcommand{\qcqs}{\mathrm{qcqs}}	% Quasi-compact quasi-seperated
\newcommand{\Qlbar}{\ov{\bQ}_\ell}			% Algebraic closure of l-adic numbers
\newcommand{\quot}{\mathrm{quot}}
\newcommand{\ra}{\rightarrow}
\newcommand{\Ra}{\Rightarrow}
\newcommand{\red}{{\mathrm{red}}}		% The maximal reduced quotient
\providecommand{\Ring}{\mathrm{Ring}}
\newcommand{\rrb}{\rrbracket}			% ]]
\newcommand{\rrp}{)\!)}			% ))
\providecommand{\Set}{\mathrm{Set}}
\newcommand{\sfp}{\mathrm{sfp}}		% strongly of finite presentation
\newcommand{\sh}{\mathrm{sh}}		% strict henselization (mainly used in superscripts)
\newcommand{\sing}{\mathrm{sing}}			% singular locus
\newcommand{\sm}{\mathrm{sm}}			% smooth locus
\providecommand{\Sn}{\text{{\upshape(S$_n$)}} }
\providecommand{\Snn}{\text{{\upshape(S$_n$)}}}
\providecommand{\Sone}{\text{{\upshape(S$_1$)}} }
\providecommand{\Sonen}{\text{{\upshape(S$_1$)}}}
\providecommand{\Stwo}{\text{{\upshape(S$_2$)}} }
\providecommand{\Stwon}{\text{{\upshape(S$_2$)}}}
\providecommand{\sqp}[1]{\left[#1\right]}
\providecommand{\sSet}{\mathrm{sSet}}
\newcommand{\st}{{\mathrm{st}}}		% semistable 
\newcommand{\surjects}{\twoheadrightarrow}
\newcommand{\syn}{\mathrm{syn}}		% for fppf cohomology (mainly used in subscripts)
\newcommand{\Tate}{\underline{\mathrm{Tate}}} 	% for Tate curve
\newcommand{\tensor}{\otimes} 			% tensor product
\newcommand{\tors}{\mathrm{tors}}		% torsion subgroup (mainly used in subscripts)
\newcommand{\un}{\underline}
\newcommand{\val}{\mathrm{val}}		% Valuation
\newcommand{\wh}{\widehat}
\newcommand{\wt}{\widetilde}
\newcommand{\xra}{\xrightarrow}
\newcommand{\Zar}{\mathrm{Zar}}		% for Zariski cohomology 

%\newcommand{\ss}{\mathrm{ss}}		% semisimplification
\newcommand{\shv}{\mathrm{shv}}		% sheaf
\renewcommand{\th}{^{\mathrm{th}}}

\newcommand{\csub[1]}{\centering\subsection{\text{#1} \nopunct} \hfill \justify}
\newcommand{\leftexp}[2]{{\vphantom{#2}}^{#1}{#2}}
\newcommand{\leftsub}[2]{{\vphantom{#2}}_{#1}{#2}}
\providecommand{\abs}[1]{\left\lvert#1\right\rvert}
\providecommand{\norm}[1]{\lVert#1\rVert}
\providecommand{\In}[1]{\left\langle#1\right\rangle}
\providecommand{\up}[1]{{\upshape(}#1{\upshape)}}
\providecommand{\resp}[1]{{\upshape(}resp.,~#1{\upshape)}}
\providecommand{\uref}[1]{{\upshape\ref{#1}}}
\providecommand{\uS}{{\upshape\S}}
\providecommand{\f}[2]{\frac{#1}{#2}}
\newcommand{\Ell}{\cE\ell\ell}				% The stack of elliptic curves
\newcommand{\EEl}{\overline{\cE\ell\ell}}		% The compactified stack of elliptic curves
\newcommand{\new}{\mathrm{new}}			% new quotient of a modular Jacobian
\newcommand{\cusps}{\mathrm{cusps}}
\newcommand{\rr}{\raggedright}
\newcommand{\bconjt}{\begin{conj-tweak}}
\newcommand{\econjt}{\end{conj-tweak}}
\newcommand{\forg}{\mathrm{forg}}




%------------------------------------------Tex shortcuts
\renewcommand{\b}{\textbf}
\providecommand{\ucolon}{{\upshape:} }
\providecommand{\uscolon}{{\upshape;} }
\providecommand{\uc}{{\upshape:} }
\providecommand{\us}{{\upshape;} }

\newcommand{\brems}{\begin{rems} \hfill \begin{enumerate}[label=\b{\thenumberingbase.},ref=\thenumberingbase]}
\newcommand{\remi}{\addtocounter{numberingbase}{1} \item}
\newcommand{\erems}{\end{enumerate} \end{rems}}
\newcommand{\begs}{\begin{egs} \hfill \begin{enumerate}[label=\b{\thenumberingbase.},ref=\thenumberingbase]}
\newcommand{\egi}{\addtocounter{numberingbase}{1} \item}
\newcommand{\eegs}{\end{enumerate} \end{egs}}
\newcommand{\m}{\item}
\newcommand{\bsm}{\begin{smallmatrix}}
\newcommand{\esm}{\end{smallmatrix}}
\newcommand{\blem}{\begin{lemma}}
\newcommand{\elem}{\end{lemma}}
\newcommand{\blemt}{\begin{lemma-tweak}}
\newcommand{\elemt}{\end{lemma-tweak}}
\newcommand{\bconj}{\begin{conj}}
\newcommand{\econj}{\end{conj}}
\newcommand{\bprob}{\begin{Problem}}
\newcommand{\eprob}{\end{Problem}}
\newcommand{\bpropt}{\begin{prop-tweak}}
\newcommand{\epropt}{\end{prop-tweak}}
\newcommand{\bq}{\begin{Q}}
\newcommand{\eq}{\end{Q}}
\newcommand{\benum}{\begin{enumerate}[label={{\upshape(\alph*)}}]}
\newcommand{\benuma}{\begin{enumerate}[label={{\upshape(\arabic*)}}]}
\newcommand{\benumr}{\begin{enumerate}[label={{\upshape(\roman*)}}]}
\newcommand{\eenum}{\end{enumerate}}
\newcommand{\bitem}{\begin{itemize}}
\newcommand{\eitem}{\end{itemize}}
\newcommand{\bc}{\begin{comment}}
\newcommand{\ec}{\end{comment}}
\newcommand{\bd}{\begin{defn}}
\newcommand{\ed}{\end{defn}}
\newcommand{\bdt}{\begin{defn-tweak}}
\newcommand{\edt}{\end{defn-tweak}}
\newcommand{\beg}{\begin{eg}}
\newcommand{\eeg}{\end{eg}}
\newcommand{\begt}{\begin{eg-tweak}}
\newcommand{\eegt}{\end{eg-tweak}}
\newcommand{\bcl}{\begin{claim}}
\newcommand{\ecl}{\end{claim}}
\newcommand{\lab}{\label}
\newcommand{\ssk}{\smallskip}
\newcommand{\msk}{\medskip}
\newcommand{\bsk}{\bigskip}
\newcommand{\x}{\text}
\newcommand{\rv}{\revise}
\newcommand{\rd}{\ready}
\providecommand{\qxq}[1]{\quad\text{#1}\quad}
\providecommand{\qx}[1]{\quad\text{#1}}
\providecommand{\xq}[1]{\text{#1}\quad}
\newcommand{\q}{\quad}
\newcommand{\qq}{\quad\quad}
\newcommand{\qqq}{\quad\quad\quad}
\newcommand{\qqqq}{\quad\quad\quad\quad}
\newcommand{\qqqqq}{\quad\quad\quad\quad\quad}
\newcommand{\qqqqqq}{\quad\quad\quad\quad\quad\quad}
\newcommand{\qqqqqqq}{\quad\quad\quad\quad\quad\quad\quad}
\newcommand{\qqqqqqqq}{\quad\quad\quad\quad\quad\quad\quad\quad}
\newcommand{\qqqqqqqqq}{\quad\quad\quad\quad\quad\quad\quad\quad\quad}
\newcommand{\tst}{\textstyle}
\newcommand{\ba}{\begin{aligned}}
\newcommand{\ea}{\end{aligned}}
\newcommand{\be}{\begin{equation}}
\newcommand{\ee}{\end{equation}}
\newcommand{\bpf}{\begin{proof}}
\newcommand{\epf}{\end{proof}}
\newcommand{\bthm}{\begin{thm}}
\newcommand{\ethm}{\end{thm}}
\newcommand{\bprop}{\begin{prop}}
\newcommand{\eprop}{\end{prop}}
\newcommand{\bcor}{\begin{cor}}
\newcommand{\ecor}{\end{cor}}
\newcommand{\brem}{\begin{rem}}
\newcommand{\erem}{\end{rem}}
\newcommand*{\QED}{\hfill\ensuremath{\qed}}
\newcommand*{\QEDD}{\hfill\ensuremath{\qed\qed}}
\newcommand{\sabcd}[4]{\left(\begin{smallmatrix}#1&#2\\#3& #4\end{smallmatrix}\right)}


%-----------------------Theorems
\usepackage{aliascnt}
\newaliascnt{numberingbase}{subsubsection}
\numberwithin{equation}{numberingbase}

\newtheoremstyle{thms}{0pt}{0pt}{\itshape}{}{\bfseries}{.}{ }{}
\theoremstyle{thms}
\newtheorem{conj}[numberingbase]{Conjecture}
\newtheorem{corollary}[numberingbase]{Corollary}
\newtheorem{cor}[numberingbase]{Corollary}
\newtheorem{lemma}[numberingbase]{Lemma}
\newtheorem{sublemma}[equation]{Lemma}
\newtheorem{prop}[numberingbase]{Proposition}
\newtheorem{proposition}[numberingbase]{Proposition}
\newtheorem{Q}[numberingbase]{Question}
\newtheorem{thm}[numberingbase]{Theorem}
\newtheorem{theorem}[numberingbase]{Theorem}
\newtheorem{variant}[numberingbase]{Variant}

\newtheoremstyle{claims}{0pt}{0pt}{}{}{\itshape}{.}{ }{}
\theoremstyle{claims}
\newtheorem{claim}[equation]{Claim}

\newtheoremstyle{defs}{0pt}{0pt}{}{}{\bfseries}{.}{ }{}
\theoremstyle{defs}
\newtheorem{defn}[numberingbase]{Definition}
\newtheorem{definition}[numberingbase]{Definition}
\newtheorem{eg}[numberingbase]{Example}
\newtheorem{example}[numberingbase]{Example}
\newtheorem*{egs}{Examples}
\newtheorem{rem}[numberingbase]{Remark}
\newtheorem{remark}[numberingbase]{Remark}
\newtheorem*{rems}{Remarks}

\Crefname{claim}{Claim}{Claims}
\Crefname{sublemma}{Lemma}{Lemmas}
\Crefname{conj}{Conjecture}{Conjectures}
\Crefname{cor}{Corollary}{Corollaries}
\Crefname{defn}{Definition}{Definitions}
\Crefname{eg}{Example}{Examples}
\Crefname{prop}{Proposition}{Propositions} 
\Crefname{Q}{Question}{Questions}
\Crefname{rem}{Remark}{Remarks}
\Crefname{thm}{Theorem}{Theorems}
\Crefname{variant}{Variant}{Variants}


\theoremstyle{thms}
\newtheorem{thm-tweak}[subsection]{Theorem}
\Crefname{thm-tweak}{Theorem}{Theorems}
\newtheorem{lemma-tweak}[subsection]{Lemma}
\Crefname{lemma-tweak}{Lemma}{Lemmas}
\newtheorem{cor-tweak}[subsection]{Corollary}
\Crefname{cor-tweak}{Corollary}{Corollaries}
\newtheorem{prop-tweak}[subsection]{Proposition}
\Crefname{prop-tweak}{Proposition}{Propositions} 
\newtheorem{conj-tweak}[subsection]{Conjecture}
\Crefname{conj-tweak}{Conjecture}{Conjectures} 

\theoremstyle{defs}
\newtheorem{defn-tweak}[subsection]{Definition}
\Crefname{defn-tweak}{Definition}{Definitions}
\newtheorem{eg-tweak}[subsection]{Example}
\Crefname{eg-tweak}{Example}{Examples}
\newtheorem*{rems-tweak}{Remarks}
\newtheorem{rem-tweak}[subsection]{Remark}
\Crefname{rem-tweak}{Remark}{Remarks}

%------------------------The bpp stuff
\newtheoremstyle{subsection-tweak}
   {0pt}
   {0pt}%
   {}
   {}%
   {\bfseries}
   {}%
   {.5em}
   {\thmnumber{\@{#1}{}\@{#2}.}%
    \thmnote{~{\bfseries#3.}}}    
    
\theoremstyle{subsection-tweak}
\newtheorem{pp}[numberingbase]{}
\newcommand{\bpp}{\begin{pp}}
\newcommand{\epp}{\end{pp}}


\theoremstyle{subsection-tweak}
\newtheorem{conventions}[subsection]{}
\newtheorem*{pp-tweak}{}



%-----------------------Enumeration
\renewcommand{\labelenumi}{\b{\arabic{enumi}.}}
\renewcommand{\labelenumii}{(\alph{enumii})}

%-----------------------Page layout
%\setlength{\parindent}{0pt} % 0 pt  = indentation
%\pagenumbering{gobble}          % Uncomment to suppress page numbers


%--------------------------Table of contents customization

\makeatletter
\def\@tocline#1#2#3#4#5#6#7{
    \begingroup 
    \@ifempty{#4}{}{}

    \parindent\z@ \leftskip#3\relax \advance\leftskip\@tempdima\relax
    #5\hskip-\@tempdima
      \ifcase #1
       \or\or \hskip 2em \or \hskip 1em \else \hskip 3em \fi%
      #6\nobreak\relax
    \dotfill\hbox to\@pnumwidth{\@tocpagenum{#7}}\par
    \nobreak
    \endgroup
 }
 \def\l@section{\@tocline{1}{0pt}{1pc}{}{}}

\renewcommand{\tocsection}[3]{%
  \indentlabel{\@ifnotempty{#2}{\makebox[1.3em][l]{%
    \ignorespaces#1 \bfseries{#2}.\hfill}}}\bfseries{#3}
    \vspace{-5pt}}

\renewcommand{\tocsubsection}[3]{%
  \indentlabel{\@ifnotempty{#2}{\hspace*{-0.5em}\makebox[2.1em][l]{%
    \ignorespaces#1#2.\hfill}}}#3
    \vspace{-5pt}}
   
%\renewcommand{\tocsectionvskip}{10pt}
%\set{10pt}
%\settocsectionvskip{10pt}

\makeatother 

\makeatletter % This is for proper formatting of appendices in the table of contents
\newcommand\appendix@section[1]{%
  \refstepcounter{section}%
  \orig@section*{Appendix \@Alph\c@section. #1}%
%  \addcontentsline{toc}{section}{Appendix \@Alph\c@section. #1}%
}
\let\orig@section\section
\g@addto@macro\appendix{\let\section\appendix@section}
\makeatother


%\setlength{\textwidth}{31pc}

%\setlength{\textheight}{48pc}

%\oddsidemargin=.65in

%\evensidemargin=.65in





%-------------------------------------------------------------END OF ALL THAT CRAP!



% Checklist to do before making a public update to the manuscript: 
%   	(a) Spell check; 
%     (b) Eliminate numbering of unreferenced equations;
%     (c) Check if cited preprints have appeared and for discrepancies in numbering;
%     (d) Thank the referee in the acknowledgements (if applicable).
%     (e) Go through the manuscript and check for unnecessary italicization of references, reference numbers, parentheses, semicolons, etc. in theorem statements.
% }







\author{K\k{e}stutis \v{C}esnavi\v{c}ius}
\address{CNRS, UMR 8628, Laboratoire de Math\'{e}matiques d'Orsay, Universit\'{e} Paris-Saclay, 91405 Orsay, France}
\email{kestutis@math.u-psud.fr}
%\urladdr{http://math.berkeley.edu/~kestutis/}

\date{\today}

%-----------------------------------------------------------------------------------------------






%\subjclass[2010]{\rv{Primary ; Secondary .}}
%\keywords{\rv{...}}


%\begin{abstract} \revise{...} \end{abstract}



%\maketitle




%\hypersetup{
%    linktoc=page,     %set to all if you want both sections and subsections linked
%}
%\renewcommand*\contentsname{}
%\q\\
%\tableofcontents




%\subsection*{Acknowledgements}  \revise{ }

 

%\begin{bibdiv} \begin{biblist} 
%\bibselect{bibliography}
%\end{biblist} \end{bibdiv}





%\usepackage[notcite]{showkeys}   % for drafts to show all the labels in pdf




\newcommand{\bb}{\boldsymbol}		% for bold symbol in math mode (esp. greek letters)


% Checklist to do before making a public update to the manuscript: 
%   	(a) Spell check; 
%     (b) Eliminate numbering of unreferenced equations;
%     (c) Check if cited preprints have appeared and for discrepancies in numbering;
%     (d) Thank the referee in the acknowledgements (if applicable).
%     (e) Go through the manuscript and check for unnecessary italicization of references, reference numbers, parentheses, semicolons, etc. in theorem statements.
% }

\usepackage{epigraph}



